\section{Results}
\label{sec:results}

% Present experimental results
[Present the main experimental results with figures and tables.]

\subsection{Performance Comparison}
\label{subsec:performance}

[Compare FedQHD against baseline methods in terms of learning performance.]

% Example figure (uncomment and modify as needed):
% \begin{figure}[htbp]
%     \centering
%     \includegraphics[width=0.8\linewidth]{figures/performance_comparison.pdf}
%     \caption{Learning curves comparing FedQHD with baseline methods.}
%     \label{fig:performance}
% \end{figure}

\subsection{Communication Efficiency}
\label{subsec:communication}

[Analyze communication efficiency and overhead.]

\subsection{Scalability Analysis}
\label{subsec:scalability}

[Examine how the method scales with the number of clients or other parameters.]

\subsection{Ablation Study}
\label{subsec:ablation}

[Conduct ablation studies to understand the contribution of different components.]

% Example table (uncomment and modify as needed):
% \begin{table}[htbp]
%     \caption{Ablation study results}
%     \label{tab:ablation}
%     \centering
%     \begin{tabular}{lcc}
%         \hline
%         \textbf{Method} & \textbf{Metric 1} & \textbf{Metric 2} \\
%         \hline
%         Full Method & 0.85 & 0.90 \\
%         Without Component A & 0.78 & 0.85 \\
%         Without Component B & 0.80 & 0.88 \\
%         \hline
%     \end{tabular}
% \end{table}
